\chapter{Image Handedness: an Error Found and Corrected}
\label{app:handedness}

After version~1.0 of this report, while an independent integrator was
being written for the real-time renderer, we found that every image the
ray tracer had produced of a \emph{rotating} system was the mirror image
of the true one. No number in this report was affected — every
quantitative result is invariant under the mirror — but which side of a
rendered disk is brighter, and which side of a Kerr shadow is flattened,
were reversed relative to the stated orientation. This appendix records
what the error was, why it arose, why the validation ladder of
Chapter~\ref{ch:devhistory} did not catch it, and how it was found,
verified and corrected. Every affected figure has been regenerated with
the corrected camera; every number quoted here is produced by
\texttt{validation/handedness\_tool} and typeset from
\texttt{numbers.tex} (Appendix~\ref{app:repro}).

\section{What a pixel records}
\label{sec:hand-pixel}

A pixel looking along the unit direction $n^{\hat\imath}$ of the ZAMO
camera (Section~\ref{sec:camera}) records the photon that \emph{arrives}
from that direction. With the local energy fixed to unity, its
four-momentum is
\begin{equation}
k_{\rm rec} = e_{\hat 0} - n^{\hat\imath}\, e_{\hat\imath},
\label{eq:krec}
\end{equation}
travelling from the sky toward the camera. To find where it came from we
follow its geodesic backwards in time, which means launching the
reversed tangent
\begin{equation}
-k_{\rm rec} = -e_{\hat 0} + n^{\hat\imath}\, e_{\hat\imath}.
\label{eq:kback}
\end{equation}
Its spatial part points along $n^{\hat\imath}$, out into the scene, as
intuition expects; its time component is negative, so the constants
extracted by \texttt{make\_null\_ray} (Section~\ref{sec:makenullray})
have $E = -k_t < 0$. Carter's equations need no modification. The
potentials $R(r)$ and $\Theta(\theta)$ are quadratic in $(E, L_z)$, so
$r(\tau)$ and $\theta(\tau)$ are exactly those of the physical light ray,
while the $t$- and $\phi$-equations, linear in $(E, L_z)$, simply run
along the backward path. The received photon carries $(-E, -L_z)$, so
the disk redshift \eqref{eq:gfactor} — whose impact-parameter form
depends only on $b_\gamma = L_z/E$ — reads, in terms of the backward
constants,
\begin{equation}
g = \frac{1}{u^t\,(\Omega L_z - E)} .
\label{eq:gback}
\end{equation}

\section{The error}
\label{sec:hand-error}

Phases 1--12 launched instead
\begin{equation}
k_{\rm legacy} = e_{\hat 0} + n^{\hat\imath}\, e_{\hat\imath},
\label{eq:klegacy}
\end{equation}
with the justification, then printed in Section~\ref{sec:camera}, that
null geodesics are time-reversible, so that where a launched ray lands
is what the pixel sees. The premise is true of a geodesic as a curve,
but traversing a curve backwards reverses \emph{all four} components of
its tangent. Reversing only the spatial part, as \eqref{eq:klegacy}
does, produces a different photon: the one the camera would
\emph{emit} toward $n^{\hat\imath}$, which is the time reverse of the
photon it receives.

In a static spacetime that distinction leaves the light paths
unchanged. In Kerr it does not. The metric \eqref{eq:kerr} is invariant
under $(t, \phi) \to (-t, -\phi)$ taken \emph{together}; the
frame-dragging term $g_{t\phi}\,\mathrm{d}t\,\mathrm{d}\phi$ changes sign
under either reversal alone. Reversing time is therefore equivalent to
reversing the azimuth, $\phi \to -\phi$: a black hole spinning the other
way. For a camera at $\phi = 0$ that reflection is through the plane
containing the camera and the spin axis, which on the image is a
left--right mirror. The same reversal sets the disk orbiting backwards,
so a legacy image is a \emph{correct} image of the mirror-image system —
which is why nothing about it looked wrong.

Algebraically, the ZAMO relations of Section~\ref{sec:camera} give the
legacy photon at pixel direction $(n^{\hat r}, n^{\hat\theta},
n^{\hat\phi})$ the constants $(-E, -L_z, Q)$ of the backward tangent
\eqref{eq:kback} at the mirror direction $(n^{\hat r}, n^{\hat\theta},
-n^{\hat\phi})$. With $R$ and $\Theta$ quadratic and the $\phi$-equation
linear in $(E, L_z)$, the two rays share $r(\tau)$ and $\theta(\tau)$
exactly and have opposite $\phi(\tau)$; comparing \eqref{eq:gback} with
the legacy formula $g = 1/[u^t(E - \Omega L_z)]$, their redshifts are
identical. The legacy pixel $(x, y)$ shows exactly what is physically
at $(-x, y)$.

Without spin the shadow is symmetric and the light paths agree, but
the disk still rotates, and the error reduces to exchanging the
Doppler-bright and Doppler-dim limbs. For the reference camera — image
up along the spin axis, image right along $+\hat\phi$ — gas orbiting
prograde (counter-clockwise seen from the pole) \emph{approaches} the
camera on the image \emph{left}, and the prograde photons that skim the
hole most closely, flattening the shadow, also arrive from the left.

\section{What was and was not affected}
\label{sec:hand-scope}

Every number in this report is mirror-invariant and therefore unchanged:
the shadow and photon-ring diameters and their comparison with the EHT
(Chapters~\ref{ch:shadow} and \ref{ch:calib}), the axis ratio, the
sub-ring flux ratios and widths (Chapter~\ref{ch:ring}), the magnitude
of the Doppler asymmetry (\BeamRatio) and the redshift range, the
Lyapunov exponents, and all S-star and no-hair results, whose orbit
integrators do not use the ray tracer at all. What changed is
orientation: the rendered images
(Figs.~\ref{fig:boundary}, \ref{fig:disk}, \ref{fig:subrings} and
\ref{fig:profile}), the sign of the horizontal centroid displacement
in Fig.~\ref{fig:shapemaps}, and two statements that had followed the
renders. Table~\ref{tab:handedness} compares the two conventions run on
the same integrator.

\begin{table}
\centering
\caption{Legacy (emitted-photon) versus corrected (received-photon)
camera on the same Carter integrator; image up along the spin axis.
Limb flux ratios are bolometric sums over the image major axis
(camera at $r = 1000\,M$, $i = 80\degr$); $g$ is the disk redshift on
either side of the hole; $\xi = L_z/E$ labels the shadow's equatorial
edges for an edge-on camera.}
\label{tab:handedness}
\begin{tabular}{lcc}
\hline
 & legacy & corrected \\
\hline
brighter limb, $\astar = 0$ (ratio) & right (\HandBeamLegacyAZero) & left (\HandBeamAZero) \\
brighter limb, $\astar = 0.94$ (ratio) & right (\HandBeamLegacyAHigh) & left (\HandBeamAHigh) \\
$g$ left / right, $\astar = 0$, $r \approx \HandRHitAZero\,M$ & \HandGRightAZero{} / \HandGLeftAZero & \HandGLeftAZero{} / \HandGRightAZero \\
shadow edge $\xi$ [$M$] left / right, $\astar = 0.94$ & $\HandXiRight$ / $\HandXiLeft$ & $\HandXiLeft$ / $\HandXiRight$ \\
\hline
\end{tabular}
\end{table}

The magnitudes are identical — only the side changes. For comparison,
the analytic Bardeen edges are $\xi = \HandXiBardeenPro\,M$ (prograde)
and $\HandXiBardeenRetro\,M$ (retrograde). (The major-axis limb ratio
of the table is a different measure from the whole-image ratio
\BeamRatio{} of Chapter~\ref{ch:disk}; both are mirror-invariant in
magnitude.)

\section{Why the validation ladder did not catch it}
\label{sec:hand-why}

\begin{itemize}
\item The one check that could have, the Gate-2 lopsidedness test,
compared only the magnitudes of the two equatorial shadow extents, with
the comment that which side is which ``depends on the camera
handedness; the asymmetry itself is the physics''. Once the camera fixes
image-up along the spin axis and image-right along $+\hat\phi$, as ours
does, the side \emph{is} physics.
\item Every other gate is mirror-invariant: the Synge radius, the
membership of a boundary ray's $(\xi, \eta)$ in the Bardeen set (a
mirrored boundary ray is simply another point of the same curve), the
diameters, the axis ratio, the sub-ring ratios, and the beaming test,
which compared the brighter limb with the dimmer one.
\item The CPU--GPU comparison of Gate~7 tested two implementations of
the same convention against each other.
\item A legacy image is a physically valid image of the mirror system,
so nothing looked wrong. Even an internal inconsistency went
unnoticed: the analytic critical curves of Fig.~\ref{fig:critcurves},
drawn in Bardeen's convention $\alpha = -\xi$, had the correct
orientation all along, while the traced boundary of
Fig.~\ref{fig:boundary} was their reflection. The two shapes differ
only by that reflection, and the axis label of the former had been
written to match the renders.
\item The shortcut is exact for the light paths of a static spacetime,
where it was first validated (Gate~1); the error enters only with
rotation, of the disk or of the hole.
\end{itemize}

\section{How it was found and verified}
\label{sec:hand-found}

The error surfaced when a second, independent integrator was written
for the real-time renderer in the Cartesian Kerr--Schild coordinates of
\eqref{eq:kerrschild} (in their outgoing form, which is regular on the
past horizon that backward-traced rays reach). That integrator uses a
Hamiltonian formulation with no Carter separation and has no polar
coordinate singularity. Its camera launcher had to be derived from
scratch in the new coordinates, and writing down which photon a pixel
receives exposed the sign of the time component in \eqref{eq:klegacy}.
Four measurements then established the diagnosis:
\begin{enumerate}
\item \emph{Mirror identity.} On \HandMirrorPixels{} disk pixels at
$\astar = 0.94$, the legacy pixel $(x, y)$ and the corrected pixel
$(-x, y)$ give the same disk-hit radius and redshift and opposite
azimuths; the largest difference is $\HandMirrorMaxDiff$. The two
computations are bit-for-bit reflections of each other, so the error is
a pure mirror and nothing else.
\item \emph{Bardeen orientation.} With the corrected camera, the
edge-on shadow edges sit at $\xi = \HandXiLeft\,M$ on the image left and
$\HandXiRight\,M$ on the right, matching the analytic prograde and
retrograde edges to $\HandXiAgree\,M$ — with the prograde edge on the
side physics requires.
\item \emph{Independent integrator.} On \HandKSPixels{} disk pixels the
Kerr--Schild tracer agrees with the corrected Carter camera to
$\HandKSOffAxisDr$ in hit radius and $\HandKSOffAxisDg$ in redshift
(relative) away from the image column above the spin axis. On that
column, where Boyer--Lindquist coordinates, and with them the Carter
reference, are singular, the difference reaches $\HandKSAxisDr$. Along
these rays the Kerr--Schild integration conserves $L_z$, $Q$ and the
null condition to $\HandKSConsL$, $\HandKSConsQ$ and $\HandKSConsH$
relative.
\item \emph{An elementary check.} For a non-rotating hole, disk gas to
the right of the hole moves away from the camera. The corrected camera
measures $g = \HandGRightAZero$ there, a redshift, and
$g = \HandGLeftAZero$ on the left; the legacy camera had them the other
way round.
\end{enumerate}

\section{The correction}
\label{sec:hand-fix}

Both science renderers, \texttt{RaytraceCpu::camera\_ray} and the
double-precision GPU shader, now launch \eqref{eq:kback} and evaluate
\eqref{eq:gback}. The legacy convention survives only behind an
explicit flag (\texttt{CameraConfig::legacy\_emitted\_photon}), set by
\texttt{handedness\_tool} to produce Table~\ref{tab:handedness} and by
a regression test that pins the mirror identity. The tests that were
deliberately side-agnostic now fix the side: Gate~2 requires the
prograde (west) shadow edge to lie nearer the look axis and prograde
boundary photons ($\xi > 0$) to arrive from the image left; Gate~3
requires the left limb to be the brighter one. New tests require the
Schwarzschild disk to be redshifted on the image right, and the Carter
and Kerr--Schild cameras to agree pixel by pixel. The real-time renderer
shared the legacy convention and has been corrected in the same way.

The affected figures were regenerated, Section~\ref{sec:camera} and
the captions of Figs.~\ref{fig:critcurves} and \ref{fig:boundary} were
corrected, and Fig.~\ref{fig:boundary} now overlays the analytic
Bardeen curve at the render's inclination, so that agreement in
orientation, not only in shape, is visible on the page. The lesson is
recorded with the other bug hunts in Section~\ref{sec:bughunts}.
